\documentclass[a4paper,12pt]{article}

\usepackage[T2A]{fontenc}
\usepackage[utf8]{inputenc}
\usepackage[russian]{babel}
\usepackage{amsmath,amssymb,mathtools}
\usepackage{helvet}
\renewcommand{\familydefault}{\sfdefault}
\usepackage{icomma}
\usepackage[a4paper,margin=2.2cm,headheight=16pt]{geometry}
\usepackage{microtype}
\usepackage{tikz}
\usepackage{pgfplots}
\pgfplotsset{compat=1.18}
\usepackage{tabularx,booktabs}
\usepackage{enumitem}
\usepackage[hidelinks]{hyperref}
\usepackage{parskip}
\usepackage{fancyhdr}
\usepackage{setspace}
\usepackage{xcolor}

\definecolor{accent}{HTML}{008080}
\definecolor{darkaccent}{HTML}{004D4D}
\definecolor{textgray}{HTML}{333333}
\definecolor{softgray}{HTML}{6B7280}
\definecolor{linegray}{HTML}{D7DEDE}

\onehalfspacing
\setlength{\parindent}{0pt}
\setcounter{secnumdepth}{0}
\setlist[itemize]{leftmargin=1.6em,itemsep=0.25em,topsep=0.35em}
\setlist[enumerate]{leftmargin=1.8em,itemsep=0.7em,topsep=0.4em}

\pagestyle{fancy}
\fancyhf{}
\fancyhead[L]{\small\textcolor{softgray}{Повторение алгебры}}
\fancyhead[R]{\small\textcolor{softgray}{Тимофей Васильченко}}
\fancyfoot[C]{\small\textcolor{softgray}{\thepage}}
\renewcommand{\headrulewidth}{0.4pt}
\renewcommand{\headrule}{\hbox to\headwidth{\color{linegray}\leaders\hrule height \headrulewidth\hfill}}

\newcommand{\key}[1]{\textcolor{darkaccent}{\textbf{#1}}}
\newcommand{\ruleline}{\par\vspace{0.3em}\noindent\textcolor{linegray}{\rule{\linewidth}{0.5pt}}\par\vspace{0.35em}}

\begin{document}

\begin{titlepage}
\thispagestyle{empty}
\centering
\vspace*{2.8cm}

{\Large\textcolor{darkaccent}{\textbf{Чек-лист}}\par}
\vspace{0.8cm}
{\huge\bfseries Повторение алгебры\par}
\vspace{0.4cm}
{\Large Степени и корни, логарифмы, тригонометрия, производная\par}

\vspace{2.2cm}
{\large Лёвин Артём Александрович\par}
\vspace{0.2cm}
{\large эксклюзивно для Тимофея Васильченко\par}

\vspace{1.2cm}
{\large \today\par}

\vfill
\begin{tikzpicture}[x=1cm,y=1cm]
  \draw[accent!45,line width=1.2pt]
    (0,0) .. controls (3.2,0.8) and (6.0,-0.7) .. (9.2,0.05)
          .. controls (11.0,0.45) and (13.0,0.35) .. (15.2,0.0);
  \draw[accent!20,line width=0.7pt]
    (1.0,-0.35) .. controls (4.1,0.1) and (7.0,-0.55) .. (10.1,-0.2)
          .. controls (12.0,0.0) and (13.7,-0.15) .. (14.6,-0.05);
\end{tikzpicture}
\end{titlepage}

\section{1. Степени и корни: привести всё к одному основанию}

\key{Базовые формулы} при $a>0$:
\[
\frac{1}{a^n}=a^{-n},\qquad
\sqrt[k]{a^m}=a^{\frac{m}{k}},\qquad
(a^m)^n=a^{mn}.
\]
Если $a>0$ и $a\ne1$, то
\[
a^{u}=a^{v}\quad\Longleftrightarrow\quad u=v.
\]

\key{Алгоритм для показательного уравнения.}
\begin{itemize}
  \item Представьте обе части как степени одного основания.
  \item Приравняйте показатели.
  \item Решите получившееся алгебраическое уравнение.
  \item Проверьте ограничения, если в исходной записи были знаменатели или корни.
\end{itemize}

\ruleline
\textbf{Пример 1.}
\[
4^{x-7}=\frac1{64}=4^{-3}
\quad\Longrightarrow\quad
x-7=-3
\quad\Longrightarrow\quad
\boxed{x=4}.
\]

\textbf{Пример 2.}
\[
\frac{\sqrt[3]{2}}{8^x}=\sqrt[5]{64}.
\]
Переходим к основанию $2$:
\[
2^{\frac13-3x}=2^{\frac65}
\quad\Longrightarrow\quad
\frac13-3x=\frac65
\quad\Longrightarrow\quad
\boxed{x=-\frac{13}{45}}.
\]

\key{Контроль.} При работе со степенями особенно внимательно проверяйте знак показателя и правило
$\frac1{a^n}=a^{-n}$.

\section{2. Иррациональные уравнения: знак правой части обязателен}

Для уравнения
\[
\sqrt{f(x)}=g(x)
\]
удобна эквивалентная система
\[
\boxed{
\sqrt{f(x)}=g(x)
\quad\Longleftrightarrow\quad
\begin{cases}
 g(x)\ge0,\\
 f(x)=g^2(x).
\end{cases}}
\]
Условие $g(x)\ge0$ защищает от лишних корней после возведения в квадрат.

\textbf{Пример 1.}
\[
\sqrt{x+6}=x
\quad\Longleftrightarrow\quad
\begin{cases}
 x\ge0,\\
 x+6=x^2.
\end{cases}
\]
\[
x^2-x-6=0
\quad\Longrightarrow\quad
x=-2\ \text{или}\ x=3.
\]
Ограничение $x\ge0$ оставляет
\[
\boxed{x=3}.
\]

\textbf{Пример 2.}
\[
\sqrt{3x+4}=-x
\quad\Longleftrightarrow\quad
\begin{cases}
 -x\ge0,\\
 3x+4=x^2.
\end{cases}
\]
Отсюда $x\le0$ и
\[
x^2-3x-4=0
\quad\Longrightarrow\quad
x=4\ \text{или}\ x=-1.
\]
Подходит только
\[
\boxed{x=-1}.
\]

\key{Типичная ошибка.} После возведения в квадрат нельзя принимать все найденные корни без проверки исходного ограничения.

\section{3. Логарифмические уравнения: определение и ОДЗ}

\key{Определение логарифма:}
\[
\boxed{\log_a b=c\quad\Longleftrightarrow\quad b=a^c},
\qquad
a>0,\quad a\ne1,\quad b>0.
\]
Мнемоническая подсказка: аргумент логарифма переходит в основание, возведённое в степень, равную значению логарифма.

Полезные преобразования:
\[
\log_a u-\log_a v=\log_a\frac{u}{v},
\qquad
k=\log_a a^k,
\]
где все аргументы логарифмов положительны.

\textbf{Пример 1.}
\[
\log_8(5x+47)=3
\quad\Longrightarrow\quad
5x+47=8^3=512,
\]
поэтому
\[
\boxed{x=93}.
\]
Здесь найденное значение автоматически даёт положительный аргумент: $5\cdot93+47=512>0$.

\textbf{Пример 2.}
\[
\log_2(x+1)=\log_2 x-1.
\]
ОДЗ:
\[
x+1>0,\qquad x>0
\quad\Longrightarrow\quad x>0.
\]
Так как $1=\log_2 2$,
\[
\log_2(x+1)=\log_2\frac{x}{2}
\quad\Longrightarrow\quad
x+1=\frac{x}{2}
\quad\Longrightarrow\quad
x=-2.
\]
Число $-2$ не принадлежит ОДЗ, следовательно,
\[
\boxed{\text{решений нет}}.
\]

\section{4. Тригонометрия: двойной угол и приведение}

\key{Формула двойного угла для синуса:}
\[
\boxed{\sin 2\alpha=2\sin\alpha\cos\alpha}.
\]

\key{Три формы для косинуса двойного угла:}
\[
\boxed{
\cos2\alpha
=\cos^2\alpha-\sin^2\alpha
=2\cos^2\alpha-1
=1-2\sin^2\alpha.}
\]

Также полезно помнить
\[
\sin(-\alpha)=-\sin\alpha,
\qquad
\cos(-\alpha)=\cos\alpha,
\qquad
\sin(\alpha+2\pi k)=\sin\alpha.
\]
Табличные углы: $0$, $\frac{\pi}{6}$, $\frac{\pi}{4}$, $\frac{\pi}{3}$, $\frac{\pi}{2}$, $\pi$.

\textbf{Пример 1.}
\[
3\sin\frac{13\pi}{12}\cos\frac{13\pi}{12}
=\frac32\sin\frac{13\pi}{6}.
\]
Поскольку
\[
\frac{13\pi}{6}=2\pi+\frac{\pi}{6},
\qquad
\sin\frac{13\pi}{6}=\frac12,
\]
получаем
\[
\boxed{\frac34=0,75}.
\]

\textbf{Пример 2.}
\[
5\cos\frac{13\pi}{8}\sin\!\left(-\frac{13\pi}{8}\right)
=-\frac52\sin\frac{13\pi}{4}.
\]
Так как
\[
\frac{13\pi}{4}=3\pi+\frac{\pi}{4},
\qquad
\sin\frac{13\pi}{4}=-\frac{\sqrt2}{2},
\]
то
\[
\boxed{\frac{5\sqrt2}{4}}.
\]

\key{Контроль знака.} После приведения угла обязательно определите четверть и знак тригонометрической функции.

\section{5. Производная: читать график $y=f'(x)$}

\key{Монотонность функции:}
\[
f'(x)>0\ \Longrightarrow\ f(x)\ \text{возрастает},
\qquad
f'(x)<0\ \Longrightarrow\ f(x)\ \text{убывает}.
\]

\key{Экстремумы для дифференцируемой функции:}
\begin{itemize}
  \item если $f'(x)$ меняет знак с $+$ на $-$, то в этой точке у $f$ локальный максимум;
  \item если $f'(x)$ меняет знак с $-$ на $+$, то в этой точке у $f$ локальный минимум;
  \item равенство $f'(x_0)=0$ само по себе экстремум не гарантирует: важна смена знака производной.
\end{itemize}

\begin{center}
\begin{tikzpicture}
\begin{axis}[
  width=10.8cm,
  height=6.2cm,
  axis lines=middle,
  unit vector ratio=1 1,
  xmin=-2.2,xmax=2.2,
  ymin=-2.0,ymax=2.0,
  samples=220,
  clip=false,
  xlabel={$x$},
  ylabel={$f'(x)$},
  xtick={-2,-1,0,1,2},
  ytick={-1,0,1},
  tick label style={font=\small},
  label style={font=\small}
]
\addplot[accent,thick,domain=-2:2] {-x};
\addplot[only marks,mark=*,mark size=2.2pt,darkaccent] coordinates {(0,0)};
\node[anchor=south west,text=darkaccent,font=\small] at (axis cs:-1.9,1.25) {$f'(x)>0$};
\node[anchor=north east,text=darkaccent,font=\small] at (axis cs:1.9,-1.25) {$f'(x)<0$};
\node[anchor=south west,text=textgray,font=\small] at (axis cs:0.08,0.12) {максимум $f$};
\end{axis}
\end{tikzpicture}
\end{center}

\key{Производная через касательную.} Если касательная к графику $y=f(x)$ в точке $x_0$ имеет угол наклона $\alpha$, то
\[
\boxed{f'(x_0)=k=\tg\alpha=\frac{\Delta y}{\Delta x}}.
\]
Для возрастающей касательной $k>0$, для убывающей $k<0$.

Например,
\[
\frac{\Delta y}{\Delta x}=\frac85=1,6
\quad\Longrightarrow\quad
f'(x_0)=1,6,
\]
а при падении на $2$ единицы за горизонтальный шаг $8$
\[
f'(x_0)=-\frac28=-0,25.
\]

\key{Формулировки ЕГЭ.}
\begin{itemize}
  \item «Точка максимума» или «точка минимума» обычно означает значение $x$.
  \item «Наименьшее значение функции» означает число $f(x)$.
  \item Если функция возрастает на всём отрезке $[a;b]$, то наименьшее значение достигается при $x=a$, наибольшее --- при $x=b$.
  \item Всегда уточняйте, чей график дан: $f(x)$ или $f'(x)$.
\end{itemize}

\section{6. Быстрый контроль перед ответом}

\begin{itemize}
  \item \key{Степени:} одинаковое основание, верный знак показателя, отсутствие нулевого знаменателя.
  \item \key{Корни:} ограничение на часть без корня, проверка корней после возведения в квадрат.
  \item \key{Логарифмы:} основание $>0$, основание $\ne1$, каждый аргумент $>0$.
  \item \key{Тригонометрия:} коэффициент перед произведением, период, четверть, знак функции.
  \item \key{Производная:} знак $f'(x)$, направление смены знака, знак углового коэффициента касательной.
\end{itemize}

\clearpage
\section{Тренировочный блок}

\textbf{3 задания среднего уровня}
\begin{enumerate}[label=\textbf{\arabic*.}]
  \item Решите уравнение
  \[
  9^{x-2}=\frac1{27}.
  \]

  \item Решите уравнение
  \[
  \sqrt{2x+3}=x.
  \]

  \item Решите уравнение
  \[
  \log_3(2x+1)=4.
  \]
\end{enumerate}

\textbf{6 заданий выше среднего}
\begin{enumerate}[label=\textbf{\arabic*.},start=4]
  \item Решите уравнение
  \[
  \frac{\sqrt[3]{4}}{8^x}=\sqrt[4]{32}.
  \]

  \item Решите уравнение
  \[
  \sqrt{x-1}=3-x.
  \]

  \item Решите уравнение
  \[
  \log_2(x+3)=\log_2 x+1.
  \]

  \item Найдите значение выражения
  \[
  4\sin\frac{7\pi}{12}\cos\frac{7\pi}{12}.
  \]

  \item Найдите значение выражения
  \[
  6\cos\frac{5\pi}{8}\sin\!\left(-\frac{5\pi}{8}\right).
  \]

  \item Известно, что
  \[
  f'(x)>0\ \text{на}\ (-5;-1)\cup(3;7),
  \qquad
  f'(x)<0\ \text{на}\ (-1;3),
  \]
  причём $f'(-1)=f'(3)=0$. Укажите точку локального максимума, точку локального минимума и промежутки возрастания функции $f$.
\end{enumerate}

\textbf{1 сложное задание}
\begin{enumerate}[label=\textbf{\arabic*.},start=10]
  \item На отрезке $[-5;4]$ производная функции задаётся формулой
  \[
  f'(x)=(x+4)(x-1)^2(x-3).
  \]
  Укажите точки локального максимума и локального минимума функции $f$, определите, является ли $x=1$ точкой экстремума, и запишите промежутки возрастания функции внутри данного отрезка.
\end{enumerate}

\clearpage
\section{Ответы к тренировочному блоку}

\renewcommand{\arraystretch}{1.35}
\begin{table}[ht]
\centering
\caption{Краткие ответы}
\begin{tabularx}{\linewidth}{@{}cX@{}}
\toprule
\textbf{№} & \textbf{Ответ} \\
\midrule
1 & $x=\dfrac12$. \\
2 & $x=3$. \\
3 & $x=40$. \\
4 & $x=-\dfrac{7}{36}$. \\
5 & $x=2$. \\
6 & $x=3$. \\
7 & $-1$. \\
8 & $\dfrac{3\sqrt2}{2}$. \\
9 & Локальный максимум: $x=-1$; локальный минимум: $x=3$; возрастание на $(-5;-1)\cup(3;7)$. \\
10 & Локальный максимум: $x=-4$; локальный минимум: $x=3$; $x=1$ экстремумом не является; возрастание на $(-5;-4)\cup(3;4)$. \\
\bottomrule
\end{tabularx}
\end{table}

\end{document}